% =====================================================================
%  FICHE 8 — THÈME 5 — NIVEAU MOYEN — ÉNONCÉS
% =====================================================================
\documentclass[11pt,a4paper]{article}
\usepackage[utf8]{inputenc}
\usepackage[T1]{fontenc}
\usepackage{lmodern}
\usepackage[french]{babel}
\usepackage{amsmath,amssymb,mathtools}
\usepackage{icomma}
\usepackage{siunitx}
\sisetup{output-decimal-marker={,},per-mode=power,group-separator={\,},group-minimum-digits=5}
\usepackage[version=4]{mhchem}
\usepackage{xcolor}
\usepackage{tikz}
\usepackage[most]{tcolorbox}
\usepackage{enumitem}
\usepackage{array}
\usepackage{fancyhdr}
\usepackage[hidelinks]{hyperref}
\usetikzlibrary{arrows.meta,positioning,calc}
\usepackage[a4paper,margin=2.1cm,headheight=26pt,headsep=10pt]{geometry}
\usepackage{titlesec}

\definecolor{primary}{RGB}{150,98,162}
\definecolor{primarydark}{RGB}{92,52,110}

\tcbset{boxbase/.style={enhanced,breakable,boxrule=0.9pt,arc=2.5pt,left=3mm,right=3mm,top=2.3mm,bottom=2.3mm,fonttitle=\bfseries,coltitle=white,toptitle=0.6mm,bottomtitle=0.6mm}}
\newtcolorbox[auto counter]{exercice}[1][]{boxbase,colback=primary!4,colframe=primary,title={\textsc{Exercice}~\thetcbcounter\ \ifx\relax#1\relax\relax\else\textmd{--- \itshape #1}\fi}}

\pagestyle{fancy}
\fancyhf{}
\fancyhead[L]{\small\textcolor{primarydark}{\textbf{Fiche 8} — Thème 5 : Équilibre chimique et Le Chatelier}}
\fancyhead[R]{\small\textcolor{primarydark}{\textsc{Moyen}}}
\fancyfoot[C]{\small\thepage}
\renewcommand{\headrulewidth}{0.8pt}
\renewcommand{\headrule}{\hbox to\headwidth{\color{primary}\leaders\hrule height \headrulewidth\hfill}}
\setlength{\parindent}{0pt}\setlength{\parskip}{3pt}

\begin{document}

\begin{center}
\begin{tikzpicture}
\node[rectangle,rounded corners=7pt,fill=primary,text=white,minimum width=16cm,
      minimum height=1.1cm,align=center,inner sep=6pt]
  {{\large\bfseries Thème 5 — Équilibre chimique et Le Chatelier}\\[1pt]{\small Niveau \textsc{Moyen} — énoncés}};
\end{tikzpicture}
\end{center}

\begin{exercice}[calculer $K$ à partir des concentrations]
Le système \;$\ce{PCl5}(g) \rightleftharpoons \ce{PCl3}(g) + \ce{Cl2}(g)$\; est à l'équilibre à
\SI{298}{\kelvin}. Les concentrations valent $[\ce{PCl5}]=\SI{0.40}{\mol\per\litre}$,
$[\ce{PCl3}]=\SI{1.40}{\mol\per\litre}$ et $[\ce{Cl2}]=\SI{0.50}{\mol\per\litre}$.
\begin{enumerate}[label=\alph*),leftmargin=2em,topsep=2pt,itemsep=2pt]
  \item Écrire l'expression de $K$.
  \item Calculer sa valeur.
  \item L'équilibre est-il plutôt du côté des réactifs ou des produits ?
\end{enumerate}
\end{exercice}

\begin{exercice}[équilibre hétérogène en phase gazeuse]
À \SI{1000}{\celsius}, la réaction \;$\ce{FeO}(s) + \ce{CO}(g) \rightleftharpoons \ce{Fe}(s) + \ce{CO2}(g)$\;
a une constante $K = 0{,}4$. On introduit du $\ce{FeO}(s)$ \emph{en excès} et du monoxyde de carbone sous
$p_{\ce{CO}} = \SI{1}{atm}$.
\begin{enumerate}[label=\alph*),leftmargin=2em,topsep=2pt,itemsep=2pt]
  \item Écrire $K$ en fonction des pressions partielles.
  \item Chaque $\ce{CO}$ consommé donne un $\ce{CO2}$ : justifier que $p_{\ce{CO}} + p_{\ce{CO2}} = \SI{1}{atm}$.
  \item En posant $p_{\ce{CO2}} = x$, calculer la pression de $\ce{CO2}$ à l'équilibre.
\end{enumerate}
\end{exercice}

\begin{exercice}[réécrire l'équation change $K$]
Pour \;$\ce{PCl5}(g) \rightleftharpoons \ce{PCl3}(g) + \ce{Cl2}(g)$,\; on mesure $K = 2$ à \SI{300}{\kelvin}.
Donner la constante d'équilibre pour :
\begin{enumerate}[label=\alph*),leftmargin=2em,topsep=2pt,itemsep=2pt]
  \item la réaction \emph{inverse} $\;\ce{PCl3}(g) + \ce{Cl2}(g) \rightleftharpoons \ce{PCl5}(g)$ ;
  \item l'équation \emph{triplée} $\;3\,\ce{PCl5}(g) \rightleftharpoons 3\,\ce{PCl3}(g) + 3\,\ce{Cl2}(g)$ ;
  \item l'équation \emph{divisée par deux} $\;\tfrac12\,\ce{PCl5}(g) \rightleftharpoons \tfrac12\,\ce{PCl3}(g) + \tfrac12\,\ce{Cl2}(g)$.
\end{enumerate}
\end{exercice}

\begin{exercice}[procédé de contact — améliorer le rendement]
À volume constant, l'oxydation \textbf{exothermique} \;$2\,\ce{SO2}(g) + \ce{O2}(g) \rightleftharpoons 2\,\ce{SO3}(g)$\;
est à l'équilibre. Pour chaque opération, dire si elle \textbf{augmente le rendement en $\ce{SO3}$}
(déplacement vers la droite) et justifier.
\begin{enumerate}[label=\alph*),leftmargin=2em,topsep=2pt,itemsep=2pt]
  \item Introduire du dioxygène en excès.
  \item Abaisser la température.
  \item Augmenter la pression.
  \item Introduire un catalyseur.
  \item Ajouter de l'argon (gaz inerte) à volume constant.
\end{enumerate}
\end{exercice}

\begin{exercice}[décomposition endothermique en récipient fermé]
Dans un récipient fermé à volume constant, la décomposition \textbf{endothermique}
\;$\ce{BaCO3}(s) \rightleftharpoons \ce{BaO}(s) + \ce{CO2}(g)$\; est à l'équilibre.
\begin{enumerate}[label=\alph*),leftmargin=2em,topsep=2pt,itemsep=2pt]
  \item Pour chaque action, dire si l'équilibre se déplace vers la \emph{droite} (formation de produits) :
        (i) augmenter $T$ ; (ii) ouvrir le récipient (le $\ce{CO2}$ s'échappe) ; (iii) ajouter du
        $\ce{BaCO3}(s)$ ; (iv) augmenter la pression ; (v) ajouter du $\ce{CO2}$.
  \item Comment évolue la valeur de $K$ lorsqu'on augmente $T$ ?
\end{enumerate}
\end{exercice}

\end{document}
